%% ==================== 问题分析 模板 (v1.5.6, BZD 数模社 2026 规范) ====================
%%
%% 借鉴 BZD 数模社 第二章 模板: 5 段连贯 + 整体求解框架图
%% 注意: BZD 建议不要再机械分 "2.1 问题一分析、2.2 问题二分析" 等孤立小节,
%%       而是用 5 段自然段体现问题联动.
%%
%% 第 1 段 (总体分析): 概括任务 + 问题关系 (递进/共享数据/参数传递)
%% 第 2 段 (数据分析): 数据来源 + 质量问题 + 预处理方法 + 选择理由
%% 第 3 段 (前序问题建模思路): 第 1 问的数学类型 + 模型选择 + 备选方案
%% 第 4 段 (后续问题承接): 后问的预测/优化/评价/情景分析
%% 第 5 段 (检验与输出): 检验方法 + 最终输出形式
%% 末尾: 整体求解框架图 (用 TikZ 或 Mermaid 跨问题流程图)
%%
%% 跑题用户必须:
%% 1. 删掉所有 % 注释 (整段 L1-L25)
%% 2. 替换【】占位符为自己的实际内容
%% 3. 删 下面的 \ProblemAnalysisSlot{...} 行 (check 15 抓这个 inline 【】)

\section{问题分析}

\textbf{【第 1 段: 总体分析, 1 段】} 题目围绕【研究对象, 1 句】提出了【概括全部任务, 1 句】。从数学建模角度看, 各问题分别涉及【评价/预测/分类/优化/决策等类型】, 并通过【数据/指标/参数/中间结果】形成较强的逻辑联系。前序问题主要用于完成【基础任务】, 所得【结果类型】将为后续【任务】提供输入; 后续问题则进一步考虑【新增条件】, 最终形成由【起点任务】到【终点任务】的完整分析过程。

\textbf{【第 2 段: 数据分析, 1 段】} 题目提供的数据主要包括【数据类型, 1 句】, 同时需要补充收集【外部数据, 1 句】。考虑到数据可能存在【缺失/异常/量纲不同/统计口径不一致/指标相关性较强】等问题, 需要先通过【数据处理方法, 1 句】完成数据清洗与口径统一, 再采用【标准化/相关性分析/主成分分析/特征筛选方法】提取有效信息, 为后续建模提供可靠数据基础。

\textbf{【第 3 段: 前序问题建模思路, 1 段】} 对于前序问题, 需要解决【任务概括, 1 句】, 其本质属于【数学问题类型, 如 评价/预测/分类/优化】。根据【数据特点和题目条件】, 可采用【主模型, 如 多元线性回归/随机森林/灰色预测/层次分析/差分进化等】刻画【变量之间的关系】, 并将得到的【权重/参数/关键指标/分类结果】传递至后续问题。为降低单一模型的局限, 可以选择【备选方法 1】和【备选方法 2】, 并依据【评价标准, 如 MAE/RMSE/AUC/约束满足率】进行比较。

\textbf{【第 4 段: 后续问题承接与扩展, 1 段】} 在前述分析的基础上, 后续问题需要进一步完成【预测/优化/影响评价/情景分析】。针对【预测任务】, 可选取【候选方法】分别刻画线性、非线性或时序变化规律, 并依据【误差指标】比较模型表现或构建组合预测; 针对【优化或决策任务】, 可将前述预测值或评价结果作为模型参数, 以【优化目标】建立【规划模型】, 同时设置【主要约束】; 对于【外部政策/环境变化/不确定性】, 可通过参数调整和情景设置分析其影响。

\textbf{【第 5 段: 检验与输出, 1 段】} 为检验模型的合理性, 可采用【误差分析/交叉验证/对比实验/灵敏度分析/稳健性检验】评价模型表现。最终输出【题目要求的预测值/评价结果/优化方案/建议】, 总体分析与求解过程如图 \ref{fig:overall_route} 所示。

\begin{figure}[htbp]
  \centering
  \begin{tikzpicture}[
    node distance=1.0cm and 1.2cm,
    box/.style={rectangle, draw, minimum width=2.8cm, minimum height=0.8cm, align=center},
    arrow/.style={->, >=stealth, thick}
  ]
    \node[box] (data) {数据预处理};
    \node[box, below=of data] (q1) {问题1: 因素分析};
    \node[box, below=of q1] (q2) {问题2: 时点优化};
    \node[box, below=of q2] (q3) {问题3: 分层优化};
    \node[box, below=of q3] (q4) {问题4: 集成决策};
    \node[box, below=of q4, fill=green!10] (check) {检验 + 灵敏度};
    \draw[arrow] (data) -- (q1);
    \draw[arrow] (q1) -- node[right]{指标/参数} (q2);
    \draw[arrow] (q2) -- node[right]{预测值} (q3);
    \draw[arrow] (q3) -- node[right]{分组结果} (q4);
    \draw[arrow] (q4) -- (check);
  \end{tikzpicture}
  \caption{本文整体求解框架}
  \label{fig:overall_route}
\end{figure}

% 跑题后必须删下面这一行 (check 15 占位符自检)
\textbf{【这里粘贴从第 1 段到第 5 段拼起来的不超过 2 页内容, 删掉这一行】}
